% !TeX program = pdflatex
\documentclass[10pt,a4paper]{article}

\usepackage[utf8]{inputenc}
\usepackage[T1]{fontenc}
\usepackage{geometry}
\usepackage[hidelinks]{hyperref}
\usepackage{xcolor}

\geometry{
  top=1.7cm,
  bottom=1.8cm,
  left=1.8cm,
  right=1.8cm
}

\pagestyle{empty}
\frenchspacing
\setlength{\parindent}{0pt}
\setlength{\parskip}{0.45em}

\definecolor{mydarkgreen}{RGB}{0,0,0}

\begin{document}

\begin{minipage}[t]{0.50\textwidth}
\textbf{\large Javier Andres TARAZONA JIMENEZ}\\
Étudiant ingénieur IA \\
Machine Learning, Vision 2D et 3D, LLMs et prototypage \\
Massy, Île-de-France, France\\
\href{mailto:javitar06@gmail.com}{javitar06@gmail.com}\\
+33 07 46 36 97 75\\
\href{https://www.linkedin.com/in/javier-andres-tarazona-jimenez-84b489222/}{LinkedIn}
\end{minipage}
\hfill
\begin{minipage}[t]{0.44\textwidth}
\raggedleft
\textbf{ENSTA}\\
Institut Polytechnique de Paris\\
Direction de la formation\\
\vspace{0.35em}
2 rue François Verny\\
29806 Brest Cedex 9\\
France\\
\vspace{0.35em}
Tél. : 02 98 34 89 88
\end{minipage}

\vspace{1.1cm}

\raggedleft
Palaiseau, le 22 juin 2026

\vspace{0.7cm}

\raggedright
\textbf{Objet : Motivation pour une troisième année en alternance à l'ENSTA}

\vspace{0.5cm}

{\small

Madame, Monsieur,

Je suis étudiant colombien en double diplôme entre l’Universidad Nacional de Colombia et l'ENSTA, au sein de l’Institut Polytechnique de Paris, avec une orientation en intelligence artificielle. Mon parcours s’est construit autour d’un intérêt profond pour les systèmes, l’abstraction mathématique et la construction de solutions concrètes. J’ai choisi l’informatique parce qu’elle permet de passer d’une idée abstraite à un système fonctionnel, maîtrisable et utile. Avec le temps, l’intelligence artificielle s’est imposée pour moi comme une évolution naturelle : non plus seulement programmer des systèmes, mais concevoir des systèmes capables de percevoir, d’apprendre, de s’adapter et d’aider à la décision.

Cette orientation correspond aussi à ma personnalité intellectuelle. Je suis attiré par les problèmes ouverts, où il n’existe pas toujours une seule réponse correcte, mais où il faut formuler des hypothèses, expérimenter, comparer des approches et prendre des décisions techniques argumentées. La vision par ordinateur représente particulièrement bien cette ambition : elle transforme les mathématiques, les données et le code en perception du monde réel. Elle relie donc mon goût pour la rigueur scientifique à une volonté d’impact concret sur l’industrie, la société et l’économie.

C’est dans cette logique que je souhaite réaliser ma troisième année à l'ENSTA en alternance. Je ne recherche pas simplement une expérience professionnelle supplémentaire, mais une mission exigeante capable de me faire passer du profil d’étudiant ayant déjà une expérience en IA à celui d’ingénieur IA capable d’intervenir dans un contexte industriel réel. Après mon expérience de recherche à l'ENSTA – U2IS en vision 3D auto-supervisée, génération de données synthétiques avec Blender et pré-entraînement de modèles avec PyTorch et Vision Transformers, je souhaite désormais confronter mes compétences à des contraintes opérationnelles : données imparfaites, performance, robustesse, validation, sécurité, utilisateurs, coût et valeur économique.

L’alternance représente pour moi le meilleur cadre pour cette transition. Elle permet d’articuler la profondeur académique de l'ENSTA avec les exigences d’une entreprise technologique. Je recherche une mission où l’intelligence artificielle est utilisée pour résoudre des problèmes réels : vision par ordinateur, modèles 3D, données industrielles, automatisation, simulation, systèmes intelligents, IA embarquée ou IA générative appliquée avec des contraintes concrètes. Une telle expérience me permettrait de renforcer mon expertise technique, d’acquérir une vision produit, de comprendre les cycles de décision industriels et de développer une crédibilité professionnelle en France.

L’entreprise d’accueil m’apporterait ainsi plusieurs éléments essentiels : une mission IA techniquement profonde, une application industrielle réelle, un environnement reconnu ou stratégique, un encadrement sérieux, des systèmes et données concrets, un contexte international et, idéalement, une continuité professionnelle après le diplôme.

En retour, je pense pouvoir apporter à l’entreprise un profil hybride : une formation solide en informatique, machine learning, statistiques et reconnaissance d’images ; une expérience pratique en vision par ordinateur, OpenCV, PyTorch, Python, Blender, architectures logicielles et pipelines de données ; ainsi qu’une capacité à passer de la recherche à l’implémentation. Chez Engin.AI, j’ai travaillé sur l’analyse d’images routières, l’optimisation de traitements OpenCV et la structuration logicielle de bibliothèques Python. À APEDYS91, j’ai contribué à une application locale de correction de texte combinant règles, LLMs locaux et contrôles de robustesse. Dans le projet Tameo, je travaille sur la vision embarquée pour un système autonome. Ces expériences m’ont appris à construire, tester, améliorer et documenter des solutions techniques dans des contextes différents.

Enfin, ce projet me semble pleinement aligné avec l’esprit de la formation d’ingénieur de l'ENSTA : former des ingénieurs capables de mobiliser des connaissances scientifiques solides, de s’adapter à des environnements complexes, de transformer la théorie en savoir-faire, de travailler avec des équipes internationales et de contribuer à des projets technologiques de grande envergure. Réaliser ma troisième année en alternance serait donc, pour moi, une manière cohérente et ambitieuse de finaliser mon parcours d’ingénieur à l'ENSTA, tout en construisant une trajectoire professionnelle claire dans l’intelligence artificielle appliquée à l’industrie.

\vspace{0.8cm}

Cordialement,

\textbf{Javier Andres TARAZONA JIMENEZ}

}

\end{document}
